% Fragmento público generado desde propietarios íntegros.
% Residencia: 52.4.2.
% Fuente: ../incoming/radion_monografia_20260827/manuscrito/sections/05_geometria_accion.tex:218-290
\paragraph{Correspondencia entre fase, acción, área y orientación}

\begin{figure}[H]
\centering
\begin{tikzpicture}[scale=1.12]
  \def\aa{4.0}
  \def\bb{2.95}
  \def\cc{2.70}
  \draw[->,RadGray] (-4.7,0)--(4.7,0) node[right]{\(Q\)};
  \draw[->,RadGray] (0,-3.5)--(0,3.5) node[above]{\(P\)};
  \fill[RadTeal!8] (0,0) ellipse[x radius=\aa,y radius=\bb];
  \draw[RadTeal,line width=1.35pt] (0,0) ellipse[x radius=\aa,y radius=\bb];
  \fill[RadCopper] (-\cc,0) circle (2pt) node[below=2mm]{\(F_-\)};
  \fill[RadCopper] (\cc,0) circle (2pt) node[below=2mm]{\(F_+\)};
  \draw[RadGold,line width=1.1pt,-{Latex}] (\aa,0)
     arc[start angle=0,end angle=57.3,x radius=\aa,y radius=\bb];
  \node[RadGold] at (3.7,2.3) {un radión};
  \draw[dashed,RadBlue] (0,0)--(\aa,0) node[midway,below]{\(a_S\)};
  \draw[dashed,RadBlue] (0,0)--(0,\bb) node[midway,left]{\(b_S\)};
  \node[align=center,RadNavy] at (0,-3.25)
    {área total \(\pi a_Sb_S=2\pi\hret=\hfull\)};
\end{tikzpicture}
\caption{La elipse de acción. Un arco paramétrico de un radián barre
\(\hret\); una vuelta encierra \(\hfull\). Los focos codifican la apertura
\(C^*\) en la carta distinguida.}
\label{fig:elipse-accion}
\end{figure}

\paragraph{Relación entre la acción reducida, la incertidumbre y el área}

La relación \(\hbar=h/(2\pi)\) suele presentarse como una conveniencia
notacional: al describir oscilaciones, rotaciones y transformadas de Fourier,
los factores \(2\pi\) desaparecen si se usa \(\hbar\). La elipse de acción
revela una lectura más fuerte. La división por \(2\pi\) es exactamente el
paso de la acción de una vuelta a la acción por unidad de fase:
\[
h=\mathcal A(2\pi),\qquad \hbar=\mathcal A(1).
\]
La constante reducida no es sólo \(h\) con una notación más cómoda; es la
densidad areal de acción respecto a la fase.

Heisenberg, Dirac y Jordan establecieron el lenguaje de variables conjugadas,
conmutadores y transformaciones canónicas. No se les atribuye aquí una
\enquote{elipse ontológica} literal sin fuente. La contribución HMT--MD es
identificar la figura positiva, construir sus semiejes desde \((A,C^*)\) y
demostrar \eqref{eq:area-theorem}. La genealogía histórica abre el lenguaje;
el teorema actual fija el objeto.

\begin{narrativebox}
La acción no se añade después al movimiento. En la elipse del radión, la
acción es el área que el movimiento genera. Una vuelta no transporta una
constante prefabricada: la encierra.
\end{narrativebox}

\Needspace{10\baselineskip}
\paragraph{Controles algebraicos de los invariantes focales}

Las semejanzas decimales se someten a control:
\begin{center}
\begin{tabularx}{.96\textwidth}{l r X}
\toprule
Comparación & Residuo & Dictamen\\
\midrule
\(d_1\) frente a \(0.54\) & \(4.5455\times10^{-3}\) & no igualdad; el \(54\) no se deduce;\\
\(d_1^2\) frente a \(8/27\) & \(2.3352\times10^{-4}\) & proximidad sin mapa;\\
\(\eta\) frente a \(3/10\) & \(-3.1032\times10^{-4}\) & no igualdad;\\
\(e_f\) frente a \(2/3\) & \(4.7852\times10^{-3}\) & \(2/3\) pertenece exactamente a \(B_0\);\\
\(e^{-\eta}\) frente a \(20/27\) & \(3.0740\times10^{-4}\) & no igualdad.\\
\bottomrule
\end{tabularx}
\end{center}
Los invariantes sólidos son \eqref{eq:eta}, \eqref{eq:excentricidad} y
\eqref{eq:invariante-focal-relativo}. El criterio HMT no consiste en premiar
un parecido decimal, sino en exigir objeto, operación y conservación.
